% ID: SEL-D0051-OA-E04
% ESTADO: cotejada. Segunda pasada visual y borrador compilado revisados.
% Fuente: Selectividad/sel_2013_matematicas_aplicadas/septiembre/reserva a/EXAMEN 3 -ANDALUCÍA.pdf
% SHA-256: 35f960392892cf4aa4b4a3e5b0a7b6150b323d584bb033f75db6a182b19206d4
% Página: 1; opción A; ejercicio original 4.

El peso de los sobres de café que fabrica una empresa sigue una ley Normal de media desconocida y desviación típica 0.3 g. Se quiere construir un intervalo de confianza para estimar dicha media, con un nivel de confianza del 98\%, y para ello se toma una muestra de 9 sobres.
\begin{enumerate}[label=\alph*)]
\item \textbf{(1 punto)} \textquestiondown{}Qué amplitud tendrá dicho intervalo?
\item \textbf{(0.5 puntos)} \textquestiondown{}Cómo afectaría a dicha amplitud un aumento del tamaño de la muestra, manteniendo el mismo nivel de confianza?
\item \textbf{(1 punto)} Obtenga el intervalo de confianza sabiendo que los pesos, en gramos, de los sobres de la muestra son:
\[
7\quad7.1\quad7\quad6.93\quad7.02\quad7\quad7.01\quad6.5\quad7.1.
\]
\end{enumerate}
